% ============================================================
%  MEM|8: Memory Before Representation
%  Recoverability, Continuation, and the Geometry of
%  Persistent Structure
%
%  Nate Guimond (Flyxion)
%  Independent Research
% ============================================================

\documentclass[12pt, oneside]{book}

% ---- packages -----------------------------------------------
\usepackage[T1]{fontenc}
\usepackage[utf8]{inputenc}
%\usepackage{lmodern}
%\usepackage{microtype}
\usepackage{geometry}
\geometry{
  paper=letterpaper,
  top=1.25in, bottom=1.25in,
  inner=1.4in, outer=1.1in,
  marginparwidth=0pt
}
\usepackage{setspace}
\onehalfspacing

% math
\usepackage{amsmath, amsthm, amssymb, mathtools}
\usepackage{bm}

% figures, tables
\usepackage{graphicx}
\usepackage{booktabs}
\usepackage{array}

% hyperlinks
\usepackage[colorlinks=true,
            linkcolor=black,
            citecolor=black,
            urlcolor=black]{hyperref}
\usepackage{url}

% headers / footers
\usepackage{fancyhdr}
\pagestyle{fancy}
\fancyhf{}
\fancyhead[LE]{\small\leftmark}
\fancyhead[RO]{\small\rightmark}
\fancyfoot[C]{\thepage}
\renewcommand{\headrulewidth}{0pt}

% chapter style
\usepackage{titlesec}
\titleformat{\chapter}[display]
  {\normalfont\Large\bfseries}
  {\chaptertitlename\ \thechapter}{12pt}{\Huge}
\titlespacing*{\chapter}{0pt}{40pt}{30pt}

% theorem environments
\theoremstyle{plain}
\newtheorem{theorem}{Theorem}[chapter]
\newtheorem{proposition}[theorem]{Proposition}
\newtheorem{lemma}[theorem]{Lemma}
\newtheorem{corollary}[theorem]{Corollary}

\theoremstyle{definition}
\newtheorem{definition}[theorem]{Definition}
\newtheorem{assumption}[theorem]{Assumption}
\newtheorem{example}[theorem]{Example}
\newtheorem{remark}[theorem]{Remark}

\theoremstyle{remark}
\newtheorem*{notation}{Notation}

% custom operators
\DeclareMathOperator{\Rec}{Rec}
\DeclareMathOperator{\dist}{dist}
\DeclareMathOperator{\sgn}{sgn}
\newcommand{\calH}{\mathcal{H}}
\newcommand{\calQ}{\mathcal{Q}}
\newcommand{\calC}{\mathcal{C}}
\newcommand{\calE}{\mathcal{E}}
\newcommand{\RR}{\mathbb{R}}
\newcommand{\CC}{\mathbb{C}}
\newcommand{\NN}{\mathbb{N}}
\newcommand{\ZZ}{\mathbb{Z}}

% ============================================================
\begin{document}

% ---- title page ---------------------------------------------
\begin{titlepage}
  \centering
  \vspace*{2cm}
  {\Huge\bfseries MEM|8\par}
  \vspace{0.8cm}
  {\LARGE Memory Before Representation\par}
  \vspace{0.5cm}
  {\large Recoverability, Continuation, and the\\
         Geometry of Persistent Structure\par}
  \vspace{2cm}
  \vspace{1cm}
  {\large Flyxion\par}
  \vspace{0.3cm}
  {\normalsize Independent Research\par}
  \vspace{0.3cm}
  {\normalsize Canada\par}
  \vspace{2cm}
  {\small Draft --- \today\par}
  \vfill
\end{titlepage}

% ---- front matter -------------------------------------------
\frontmatter

% Epigraph
\clearpage
\vspace*{5cm}
\begin{center}
\itshape
A memory is not a preserved past.\\[4pt]
A memory is the continued ability\\
to regenerate distinctions from transformation.
\end{center}
\vfill
\clearpage

\tableofcontents

% ---- Preface ------------------------------------------------
\chapter*{Preface}
\addcontentsline{toc}{chapter}{Preface}

This monograph develops the MEM|8 framework from first principles.
The central claim is simple but has far-reaching consequences:
\emph{memory is not storage; memory is recoverable continuation}.

No physical system preserves a past state unchanged.
Every substrate undergoes continual transformation at every timescale.
The question is therefore never whether a state remains, but
whether the distinctions that once were present can be reconstructed
from the state that now exists.

That shift from storage to recoverability is not a minor
revision of memory theory.
It is a change of ontological category.
Traditional accounts place representation first and derive
memory as its persistence.
MEM|8 reverses the order:
recoverable distinctions are primary, and representations emerge
from stable patterns of recoverability.

The formal architecture rests on three empirical anchors
drawn from contemporary physics, neuroscience, and cell biology.
The topological acoustic synapse (TAS) provides a
physical demonstration that a small substrate can generate
a large high-dimensional space of reconstructive witnesses.
Intrinsic macroscale oscillatory modes in fMRI data provide
the dynamical picture: memory states are resonances of
a latent modal basis, and recall is thresholded excitation
rather than lookup.
Synaptic vesicle trafficking provides the maintenance picture:
persistence is not component preservation but flux balance,
a continuous renewal that conserves reconstructive organization
while replacing every material carrier.

Together these three anchors yield the MEM|8 synthesis:

\[
  \boxed{\;
    \text{Memory}
    =
    \text{recoverable resonant structure maintained by flux}.
  \;}
\]

The monograph is arranged in eight parts.
Parts~I--III build the mathematical foundations.
Parts~IV--V treat ecological and social memory.
Part~VI connects MEM|8 to computation.
Part~VII extends the theory to civilization scale.
Part~VIII addresses the foundational inversions:
representation, identity, and time all emerge from memory
rather than preceding it.

Mathematical prerequisites are a first course in real analysis,
linear algebra, and ordinary differential equations.
Familiarity with Hilbert spaces is helpful but not assumed;
the necessary functional analysis is developed inline.

\bigskip
\noindent
Flyxion\\
Canada, \today

% ---- main matter --------------------------------------------
\mainmatter

% ============================================================
\part{The Failure of Storage Metaphors}
% ============================================================

\chapter{The Storage Illusion}

\section{Why Storage Cannot Be Primitive}

The dominant picture of memory inherited from computer science
treats memory as a container: states are written in, preserved,
and read out.
This picture is useful as an engineering abstraction, but it
misrepresents the physical situation.

\begin{proposition}[No Physical Substrate Is Stably Inert]
\label{prop:no-inert}
No physical substrate under realistic conditions preserves
an exact microstate indefinitely.
More precisely, for any substrate $S$ at finite temperature
$T > 0$, the probability that $S(t) = S(t_0)$ exactly
decays to zero as $t \to \infty$.
\end{proposition}

\begin{proof}
Under Hamiltonian dynamics a substrate may return
arbitrarily close to an earlier state (Poincaré recurrence),
but exact microstate equality $S(t) = S(t_0)$ for $t > t_0$
is measure-zero in phase space.
At finite temperature, coupling to a thermal bath introduces
stochastic forcing $\eta(t)$; the Fokker--Planck equation
for the state distribution has no point-mass stationary
solution at $S(t_0)$, so the probability of exact
preservation decays strictly with time.
Under dissipative dynamics without thermal noise a fixed
point may exist, but biological and cognitive substrates
operate far from equilibrium and are therefore excluded
from this case.
In all realistic conditions, exact microstate preservation
is either measure-zero or dynamically unstable under
perturbation.
\end{proof}

\begin{remark}
The proposition is not that states never recur approximately
--- they may, as harmonic oscillators demonstrate.
The point is that exact preservation is not a stable property:
any perturbation, however small, eventually destroys it.
Memory therefore cannot be grounded in exact preservation;
it must be grounded in something more robust.
\end{remark}

The consequence is immediate: if memory required exact preservation
of a past state, no physical system could be said to remember
anything.
Since manifestly some systems do remember, the storage account
must be wrong.

\begin{definition}[Storage Account]
A \emph{storage account} of memory holds that a system $S$
remembers $x$ at time $t$ if and only if there exists a
subsystem $\sigma \subseteq S$ such that
\[
  \sigma(t) = \sigma(t_0)
  \quad\text{and}\quad
  \sigma(t_0) \text{ encodes } x.
\]
\end{definition}

Proposition~\ref{prop:no-inert} shows that the antecedent
of the storage account is never physically satisfied.
The account must therefore be abandoned or radically revised.

\section{What Survives Change}

Something clearly does survive.
A person remembers their childhood; a city retains the imprint
of its construction; a document preserves the intent of its
author.
The question is what kind of thing survives when no state
literally persists.

The answer proposed here is a \emph{capacity}:
the capacity to regenerate, from the current state, the
distinctions that were once present.

\begin{definition}[Reconstructive Capacity]
\label{def:recon-capacity}
A system $S$ has \emph{reconstructive capacity} for
a distinction $\delta$ at time $t$ if there exists a
physically realizable process
$\mathcal{R}$
such that
\[
  \mathcal{R}(S(t)) \;\approx_\varepsilon\; \delta,
\]
where $\approx_\varepsilon$ denotes approximate recovery
within tolerance $\varepsilon > 0$.
\end{definition}

This shifts the question from ``what is stored?'' to
``what can be reconstructed?''
The shift is more than terminological.
It opens the way to a theory in which memory is graded,
ecological, and dynamical rather than binary, internal,
and static.

\chapter{Memory as Reconstruction}

\section{Reconstruction Operators}

\begin{definition}[Reconstruction Operator]
Let $\calH$ be the space of admissible states of a system
and let $\calQ$ be a space of distinctions (queries, labels,
categories).
A \emph{reconstruction operator} is a map
\[
  R_t : \calH \;\longrightarrow\; \calQ,
\]
indexed by current time $t$, whose intended action is to
recover, from the present state $S(t) \in \calH$, the
distinction that was present at the target time $t_0$.
\end{definition}

\begin{definition}[Memory Strength]
\label{def:mem-strength}
Given a metric $d_{\calQ}$ on the distinction space,
the \emph{memory strength} of a system for distinction
$q_0 = R_{t_0}(S(t_0))$ at current time $t$ is
\[
  M(t_0, t)
  \;=\;
  1 - d_{\calQ}\!\bigl(R_t(S(t)),\; q_0\bigr).
\]
Perfect memory corresponds to $M = 1$; complete forgetting
to $M = 0$ (or below some operational threshold).
\end{definition}

\begin{remark}
Memory strength is not a property of the substrate alone.
It depends jointly on the current state $S(t)$,
the reconstruction operator $R_t$,
and the target distinction $q_0$.
This three-way dependence is absent from storage accounts,
which locate memory entirely in the substrate.
\end{remark}

\section{Admissible Recovery}

Not every reconstruction is admissible.
A system that produces an arbitrary output trivially
achieves ``reconstruction'' in a vacuous sense.
Admissibility constrains which outputs count.

\begin{definition}[Admissibility Metric]
\label{def:admissibility-metric}
An \emph{admissibility metric} on $\calQ$ is a metric
$d_{\calQ}$ such that two distinctions $q, q'$
are considered the same recovered memory if and only if
$d_{\calQ}(q, q') \leq \varepsilon$ for a contextually
fixed tolerance $\varepsilon$.
\end{definition}

\begin{proposition}[Memory Is a Pseudometric]
Define
\[
  \rho(t_0, t)
  = d_{\calQ}\!\bigl(R_t(S(t)),\; R_{t_0}(S(t_0))\bigr).
\]
Then $\rho$ is a pseudometric on the space of
$(t_0, t)$ pairs.
\end{proposition}

\begin{proof}
\begin{enumerate}
\item \emph{Non-negativity.} $\rho \geq 0$ because $d_{\calQ}$
      is a metric.
\item \emph{Identity.} $\rho(t_0, t_0) = 0$ because
      $R_{t_0}(S(t_0)) = R_{t_0}(S(t_0))$.
\item \emph{Symmetry.} $\rho(t_0,t) = \rho(t, t_0)$ by
      symmetry of $d_{\calQ}$.
\item \emph{Triangle inequality.}
      \begin{align*}
        \rho(t_0, t_2)
        &= d_{\calQ}(R_{t_2}(S(t_2)),\, R_{t_0}(S(t_0)))\\
        &\leq d_{\calQ}(R_{t_2}(S(t_2)),\, R_{t_1}(S(t_1)))
           + d_{\calQ}(R_{t_1}(S(t_1)),\, R_{t_0}(S(t_0)))\\
        &= \rho(t_1, t_2) + \rho(t_0, t_1).
      \end{align*}
\end{enumerate}
Therefore $\rho$ satisfies all pseudometric axioms.
(It fails to be a metric because two states with identical
reconstructive outputs may differ as physical states.)
\end{proof}

\chapter{The Epistemic Necessity of Memory}

\begin{definition}[Distinction Capacity]
Write $\delta_R(t)$ for the number of distinct outputs
a reconstruction operator $R_t$ can discriminate among
at time $t$.
\end{definition}

\begin{theorem}[Collapse of Temporal Comparison]
\label{thm:collapse}
If $\delta_R(t) = 0$ for all $t$, then no temporal
comparison, learning, or reference is possible.
\end{theorem}

\begin{proof}
Temporal comparison requires distinguishing the state
at $t_0$ from the state at $t_1$.
If $\delta_R = 0$, every reconstruction is identical,
so no pair $(t_0, t_1)$ is distinguishable by any query.
Learning requires detecting a change between an initial
state and a later state; if no states are distinguishable,
no change is detectable.
Reference requires identifying a past target; with
$\delta_R = 0$ every past target maps to the same output.
All three capacities therefore collapse to zero.
\end{proof}

\begin{corollary}
Cognition, science, and language are only possible
in systems that maintain nonzero reconstructive capacity.
Memory, in the sense of Definition~\ref{def:recon-capacity},
is a necessary condition for any epistemic activity.
\end{corollary}

% ============================================================
\part{Event Memory}
% ============================================================

\chapter{Events Before Objects}

\section{The Primacy of Events}

The storage account places objects first: there are things
that exist, and memory stores facts about them.
MEM|8 inverts this order.
Objects are not given; they are reconstructed.
What is given are events, and objects are the stable patterns
that emerge from reconstruction across event sequences.

\begin{definition}[Event Sequence]
An \emph{event sequence} is a finite or countably infinite
ordered collection
\[
  E = (e_1, e_2, \ldots, e_n, \ldots)
\]
where each $e_i$ belongs to an event space $\mathcal{E}$.
\end{definition}

\begin{definition}[Object as Compressed History]
An \emph{object identity} $\mathcal{O}$ relative to an
event sequence $E$ and reconstruction operator $R$ is
an equivalence class
\[
  \mathcal{O} = \{e_i \in E : R(H_i) \approx_\varepsilon R(H_j)
  \text{ for all } i, j\}
\]
where $H_i = \{e_k : k \leq i\}$ is the history up to
event $i$.
\end{definition}

In words: an object is whatever a reconstruction operator
assigns consistently the same answer across a stretch of
the event sequence.
Object identity is derived, not primitive.

\chapter{Histories as Computational Structures}

\begin{definition}[Historical State]
The \emph{historical state} at time $t$ is
\[
  H_t = \{e_i \mid i \leq t\},
\]
the complete record of events up to and including $t$.
\end{definition}

\begin{definition}[Compression Ratio]
Let $K(H_t)$ denote the Kolmogorov complexity of $H_t$.
The \emph{compression ratio} of the present state $S(t)$
as a witness to $H_t$ is
\[
  \chi(t) = \frac{K(H_t)}{|S(t)|},
\]
where $|S(t)|$ is the description length of the current
physical state.
A system is an \emph{efficient witness} of its history when
$\chi(t) \gg 1$.
\end{definition}

\begin{proposition}[Present State as Compressed Witness]
The present physical state of a system is a lossy compression
of its historical trajectory.
Memory is the capacity to partially decompress this compression
on demand.
\end{proposition}

\begin{proof}[Sketch]
The present state $S(t)$ is a function of all past states
via the dynamics: $S(t) = F(S(0), e_1, \ldots, e_t)$.
The map $F$ is generally not injective; many distinct
histories $H_t$ can produce the same $S(t)$.
Therefore $S(t)$ is a compression of $H_t$, and
$R_t$ is the corresponding decompression applied selectively
to the query at hand.
\end{proof}

\chapter{The Event-Log Ontology}

\begin{definition}[Event-Log System]
A system $\Sigma$ is an \emph{event-log system} if its
current state can be interpreted as a compressed summary
of an event log and if that summary supports
at least partial reconstruction of historical distinctions.
\end{definition}

\begin{example}
The following are event-log systems under this definition:
\begin{enumerate}
\item \emph{Biological cells.} The epigenetic state summarizes
      developmental and environmental event history.
\item \emph{Brains.} Synaptic weights compress a lifetime
      of sensorimotor events.
\item \emph{Documents.} Text preserves the communicative
      acts that produced it.
\item \emph{Institutions.} Policies, procedures, and
      personnel encode accumulated decisions.
\item \emph{Languages.} Phonological and grammatical structures
      reflect historical contact events across centuries.
\item \emph{Ecosystems.} Species composition records
      evolutionary and climatic history.
\end{enumerate}
\end{example}

The event-log ontology unifies these diverse systems under
a single memory-theoretic description.
Memory is not a special faculty of minds; it is a property
of any system whose present state is a compressed witness
of past events.

% ============================================================
\part{The Mathematical Architecture of Recall}
% ============================================================

\chapter{Memory as Recoverable Modal Structure}

\section{The Memory Field}

\begin{definition}[Memory Field]
\label{def:memory-field}
Let $\Omega$ be a cognitive or physical domain and let
\[
  \calH = L^2(\Omega;\CC)
\]
be the Hilbert space of admissible memory states.
A \emph{memory field} is a time-indexed section
\[
  M : \RR_{\geq 0} \;\longrightarrow\; \calH,
  \qquad t \mapsto M(t),
\]
whose instantaneous state admits a modal expansion
\[
  M(t) = \sum_{i \in I} a_i(t)\,\psi_i,
\]
where $\{\psi_i\}_{i \in I}$ is an orthonormal family
of \emph{reconstructive modes} and $a_i(t) \in \CC$
are time-dependent excitation coefficients.
\end{definition}

The modes $\{\psi_i\}$ are latent structures of the memory
system; they are determined by the physical or cognitive
substrate, not by the content of any individual memory.
Individual memories are patterns of excitation over this
fixed modal basis.

\section{Memory Is Not Substrate Identity}

\begin{definition}[Memory Equivalence]
\label{def:mem-equiv}
Given a reconstruction functional $R : \calH \to \calQ$,
two states $M, N \in \calH$ are
\emph{memory-equivalent relative to $R$} if
\[
  M \sim_R N
  \;\;\Longleftrightarrow\;\;
  d_{\calQ}(R(M),\, R(N)) \leq \varepsilon.
\]
A memory \emph{persists} over $[t_0, t_1]$ when
$M(t) \sim_R M(t_0)$ for all $t \in [t_0, t_1]$,
regardless of whether the coefficients $a_i(t)$ or
the material carriers of $M(t)$ change.
\end{definition}

\begin{proposition}[Memory Is Not Substrate Identity]
\label{prop:not-substrate}
If $M(t) = \sum_i a_i(t)\psi_i$ and
$N(t) = \sum_i b_i(t)\psi_i$ satisfy
$d_{\calQ}(R(M(t)), R(N(t))) \leq \varepsilon$,
then $M(t)$ and $N(t)$ instantiate the same recoverable
memory relative to $R$, even if
$a_i(t) \neq b_i(t)$ for every mode $i$.
\end{proposition}

\begin{proof}
Memory equivalence (Definition~\ref{def:mem-equiv}) depends
only on the image under $R$ in $\calQ$, not on equality
in $\calH$.
Since $d_{\calQ}(R(M(t)), R(N(t))) \leq \varepsilon$
holds by hypothesis, the two field states are
memory-equivalent for every query stable under tolerance
$\varepsilon$.
Therefore memory persists as a recoverable distinction
rather than as a conserved material or modal configuration.
\end{proof}

\chapter{Spectral Decomposition of Memory Fields}
\label{ch:spectral}

\section{The Memory Operator}

The modal expansion $M(t) = \sum_i a_i(t)\psi_i$
introduced in the previous chapter requires justification.
The modes $\{\psi_i\}$ do not simply appear; they are
eigenfunctions of a natural operator associated with the
memory system.
This chapter derives the modal basis from spectral theory,
so that all subsequent modal machinery rests on a firm
functional-analytic foundation.

\begin{definition}[Memory Operator]
Let $\calH = L^2(\Omega;\CC)$ be the memory field Hilbert space.
A \emph{memory operator} is a bounded, self-adjoint linear
operator
\[
  \mathcal{L} : \calH \;\longrightarrow\; \calH
\]
that encodes the intrinsic reconstructive structure of the
cognitive domain $\Omega$.
Physically, $\mathcal{L}$ captures how the system propagates
distinctions across space and frequency: high eigenvalues
correspond to modes that are strongly supported by the substrate;
low eigenvalues correspond to modes that are weakly supported
and therefore easily overdamped.
\end{definition}

\begin{theorem}[Spectral Decomposition of $\mathcal{L}$]
\label{thm:spectral}
Let $\mathcal{L}$ be a compact, self-adjoint operator on
$\calH$.
Then there exists a countable orthonormal family
$\{\psi_i\}_{i \in \NN} \subset \calH$ and a sequence of
real eigenvalues $\mu_1 \geq \mu_2 \geq \cdots \to 0$
such that
\[
  \mathcal{L}\psi_i = \mu_i\,\psi_i
  \qquad \text{for all } i \in \NN,
\]
and every $M \in \calH$ admits the expansion
\[
  M = \sum_{i=1}^\infty \langle M, \psi_i\rangle\,\psi_i,
\]
convergent in the $\calH$-norm.
\end{theorem}

\begin{proof}
By the spectral theorem for compact self-adjoint operators
on a Hilbert space (see, e.g., Reed and Simon \cite{ReedSimon1972},
Theorem VI.16), any compact self-adjoint operator admits a
countable orthonormal eigenbasis with real eigenvalues
accumulating only at zero.
The expansion of $M$ then follows from completeness of
the eigenbasis.
\end{proof}

\begin{definition}[Reconstructive Modes]
The eigenfunctions $\{\psi_i\}$ of $\mathcal{L}$ are the
\emph{reconstructive modes} of the memory system.
The eigenvalue $\mu_i$ measures the \emph{modal support}:
the degree to which the substrate amplifies or sustains
mode $\psi_i$.
Modes with large $\mu_i$ are intrinsically stable;
modes with small $\mu_i$ are weakly sustained and
vulnerable to overdamping.
\end{definition}

\section{Time Evolution in the Modal Basis}

Once the spectral basis is established, the dynamics of
the memory field are naturally expressed modally.

\begin{proposition}[Modal Coefficient Dynamics]
\label{prop:modal-dyn}
Let $M(t) = \sum_i a_i(t)\psi_i$ be the modal expansion
of the memory field, and suppose the field evolves by
\[
  \frac{\partial M}{\partial t}
  = -\mathcal{L}^\dagger M + F(t),
\]
where $\mathcal{L}^\dagger$ is a damping operator diagonal
in the modal basis with entries $(\lambda_i + i\omega_i)$,
and $F(t)$ is an external forcing term.
Then each coefficient satisfies
\[
  \frac{d a_i}{dt}
  = -(\lambda_i + i\omega_i)\,a_i + f_i(t),
\]
where $f_i(t) = \langle F(t), \psi_i\rangle$.
\end{proposition}

\begin{proof}
Project the field equation onto mode $\psi_i$ by taking
the inner product with $\psi_i$:
\[
  \frac{d}{dt}\langle M(t), \psi_i\rangle
  = -\langle \mathcal{L}^\dagger M(t), \psi_i\rangle
    + \langle F(t), \psi_i\rangle.
\]
Since $\mathcal{L}^\dagger$ is diagonal in the $\{\psi_i\}$
basis, $\langle \mathcal{L}^\dagger M, \psi_i\rangle
= (\lambda_i + i\omega_i)\langle M, \psi_i\rangle
= (\lambda_i + i\omega_i)a_i$.
Substituting $a_i(t) = \langle M(t), \psi_i\rangle$ gives
the stated ODE.
\end{proof}

This proposition retroactively justifies the
``Damped Memory Mode'' definition of Chapter~9: each
coefficient ODE is not an assumption but a consequence of
projecting the field equation onto the spectral basis
supplied by Theorem~\ref{thm:spectral}.

\section{The Modal Persistence Theorem}

The spectral framework now enables a theorem that connects
Part~III directly to Part~V.
Persistence-through-flux (Chapter~13) shows that
reconstructive density is conserved under balanced renewal.
The following theorem shows that memory persists under
arbitrary drift of modal coefficients, provided the
reconstruction functional remains stable.

\begin{theorem}[Modal Persistence]
\label{thm:modal-persistence}
Suppose $M(t) = \sum_i a_i(t)\psi_i$ and let
$R : \calH \to \calQ$ be a reconstruction functional
that is continuous with respect to the $\calH$-norm.
If
\[
  \|R(M(t)) - R(M(t_0))\|_{\calQ} < \varepsilon
\]
for all $t \in [t_0, t_1]$, then memory persists over
$[t_0, t_1]$ despite arbitrary variation of the
individual coefficients $a_i(t)$.
\end{theorem}

\begin{proof}
Memory persistence is defined by memory equivalence
(Definition~\ref{def:mem-equiv}): $M(t) \sim_R M(t_0)$
whenever $d_{\calQ}(R(M(t)), R(M(t_0))) \leq \varepsilon$.
The hypothesis states exactly this condition.
Therefore, regardless of how the coefficients $a_i(t)$
vary --- they may oscillate, drift, grow, or shrink ---
the memory persists as long as the reconstruction
functional's output remains within tolerance $\varepsilon$
of its initial value.
The coefficients are implementation details of the field;
the reconstruction output is what constitutes the memory.
\end{proof}

\begin{corollary}[Robustness of Memory to Substrate Fluctuation]
Under the conditions of Theorem~\ref{thm:modal-persistence},
a memory survives arbitrary fluctuation of the modal
coefficients $a_i(t)$ provided the fluctuations collectively
preserve the output of $R$.
In particular, memory is robust to any transformation
$a_i \mapsto a_i + \delta_i$ satisfying
\[
  \left\|R\!\left(\sum_i (a_i + \delta_i)\psi_i\right)
        - R\!\left(\sum_i a_i\psi_i\right)\right\|_{\calQ}
  < \varepsilon.
\]
\end{corollary}

\begin{proof}
Substitute the perturbed coefficients into the condition
of Theorem~\ref{thm:modal-persistence} and apply the
definition of memory equivalence.
\end{proof}

The Modal Persistence Theorem is the bridge between the
dynamical picture (modes oscillate and damp) and the
ecological picture (memories survive through flux renewal):
what matters is not the trajectory of any individual
coefficient but the collective stability of what the
reconstruction functional can recover.

\chapter{Dimensional Amplification and Witness Coordinates}

\section{Witness Expansion}

Empirical work on the topological acoustic synapse (TAS)
demonstrates that a small physical substrate driven by
$k$ input frequencies can generate an exponentially larger
family of nonlinear interaction states --- a
\emph{phi-bit lattice} --- whose effective dimensionality
far exceeds the input dimension \cite{Chen2026TAS}.
This section formalizes the general principle as the
\emph{Dimensional Amplification Theorem}.

\begin{definition}[Witness Expansion]
\label{def:witness-expansion}
Let $X$ be an input space of events, cues, or sensory
fragments.
A \emph{witness expansion} is a map
\[
  \Phi : X \;\longrightarrow\; \CC^n,
  \qquad x \mapsto (\phi_1(x), \ldots, \phi_n(x)),
\]
where each coordinate $\phi_j$ is a \emph{reconstructive
witness}.
The expansion is \emph{memory-useful} when there exists
a low-complexity readout
\[
  L : \CC^n \;\longrightarrow\; \calQ
\]
such that $L(\Phi(x)) \approx Q(x)$ for the relevant
query $Q : X \to \calQ$.
\end{definition}

\begin{definition}[Reconstructive Separability]
Two event classes $A, B \subset X$ are
\emph{reconstructively separable} by $\Phi$ when there
exists a hyperplane vector $h \in \CC^n$ such that
\[
  \operatorname{Re}\langle h, \Phi(a)\rangle > 0
  \quad\text{and}\quad
  \operatorname{Re}\langle h, \Phi(b)\rangle < 0
\]
for all $a \in A$ and $b \in B$ outside an admissible
error set of measure zero.
\end{definition}

\begin{theorem}[Dimensional Amplification Theorem]
\label{thm:dim-amp}
Let $\Phi : X \to \CC^n$ be a nonlinear witness expansion.
If two classes $A, B \subset X$ are not linearly separable
in $X$ but become linearly separable in $\Phi(X)$,
then memory retrieval may be implemented by a simpler
readout in witness space than in the original event space.
\end{theorem}

\begin{proof}
Assume $A, B \subset X$ are not linearly separable in
the original event coordinates.
Suppose that under $\Phi$ there exists $h \in \CC^n$
such that
\[
  \sgn\!\operatorname{Re}\langle h, \Phi(x)\rangle
\]
correctly distinguishes $A$ from $B$.
Then the decision boundary in $X$ is the pullback
\[
  \Phi^{-1}\{z : \operatorname{Re}\langle h, z\rangle = 0\},
\]
which may be a nonlinear hypersurface in $X$ while remaining
a hyperplane in $\Phi(X)$.
The classification complexity has therefore been absorbed
into the structure of the expansion $\Phi$ rather than
the readout $L$.
Retrieval is easier because the memory system searches
the amplified witness space, not the original event space.
\end{proof}

\begin{corollary}[Memory as Witness Redundancy]
\label{cor:witness-redundancy}
If a memory is represented by $n$ witnesses and
reconstruction requires only $k < n$ witnesses,
then the memory is robust to the loss of any $n - k$
witness coordinates, provided the remaining $k$
coordinates determine the same query answer.
\end{corollary}

\begin{proof}
Let $W = \{\phi_1, \ldots, \phi_n\}$ be the full witness
set and $W' \subseteq W$ with $|W'| = k$.
If $R_W(M) = R_{W'}(M)$ within tolerance $\varepsilon$,
then deletion of the $n - k$ unused witnesses does not
alter the recoverable distinction.
Therefore the memory persists under partial witness loss
whenever the remaining witnesses form a sufficient set
for the query.
\end{proof}

\chapter{Resonant Modal Recall}

\section{Cue Coupling and Modal Dynamics}

Functional MRI studies of intrinsic oscillatory modes
reveal that large-scale brain activity decomposes as
$\Psi(n, t) = \sum_\alpha \psi_\alpha(n)\tau_\alpha(t)$,
where a small family of spatial modes $\psi_\alpha$
accounts for the majority of observed long-range
connectivity \cite{Cabral2023OscillatoryModes}.
MEM|8 reinterprets this finding: the modes are not
incidental features of neural anatomy but the latent
basis of a reconstructive memory field.
Recall is resonance with that basis, not retrieval
from a store.

\begin{definition}[Cue Coupling]
A cue $c$ \emph{couples} to memory mode $\psi_i$ with
strength
\[
  \gamma_i(c) = |\langle C(c),\, \psi_i\rangle|,
\]
where $C(c) \in \calH$ is the field representation
of the cue.
\end{definition}

\begin{definition}[Damped Memory Mode]
A memory mode has activation $a_i(t) \in \CC$ governed by
\[
  \frac{d a_i}{dt}
  = -(\lambda_i + i\omega_i)\,a_i
    + \gamma_i(c(t))
    + \eta_i(t),
\]
where $\lambda_i \geq 0$ is the \emph{damping coefficient},
$\omega_i$ the \emph{modal frequency},
$\gamma_i(c(t))$ is cue forcing, and
$\eta_i(t)$ is background excitation (noise).
\end{definition}

\section{The Recall Threshold}

\begin{definition}[Ecphoric Threshold]
The \emph{ecphoric threshold} $\theta_i > 0$ is the
minimum activation magnitude required for mode $\psi_i$
to contribute to conscious recall:
\[
  \text{mode } \psi_i \text{ is recall-active}
  \;\;\Longleftrightarrow\;\;
  |a_i(t)| \geq \theta_i.
\]
\end{definition}

\begin{theorem}[Resonant Recall Criterion]
\label{thm:resonant-recall}
For constant cue forcing $\gamma_i$ and negligible noise,
the steady-state amplitude satisfies
\[
  |a_i^*| = \frac{|\gamma_i|}{\sqrt{\lambda_i^2 + \omega_i^2}}.
\]
Recall occurs precisely when
\[
  |\gamma_i| \;\geq\; \theta_i\sqrt{\lambda_i^2 + \omega_i^2}.
\]
\end{theorem}

\begin{proof}
For constant forcing the ODE reduces to
\[
  \frac{d a_i}{dt} = -(\lambda_i + i\omega_i)\,a_i + \gamma_i.
\]
The unique equilibrium is
\[
  a_i^* = \frac{\gamma_i}{\lambda_i + i\omega_i}.
\]
Taking magnitudes:
\[
  |a_i^*|
  = \frac{|\gamma_i|}{|\lambda_i + i\omega_i|}
  = \frac{|\gamma_i|}{\sqrt{\lambda_i^2 + \omega_i^2}}.
\]
Recall requires $|a_i^*| \geq \theta_i$, which gives
\[
  |\gamma_i| \geq \theta_i\sqrt{\lambda_i^2 + \omega_i^2},
\]
establishing the threshold condition.
\end{proof}

\begin{remark}
Recall is not lookup.
It is thresholded resonant excitation.
Two factors govern whether a cue triggers recall:
cue-mode overlap (the numerator $|\gamma_i|$)
and modal impedance (the denominator
$\sqrt{\lambda_i^2 + \omega_i^2}$).
Strong cues with high coupling always recall;
weak or mismatched cues fail, even when the mode itself
is intact.
\end{remark}

\chapter{Forgetting as Overdamping, Repair as Resonance Restoration}

\section{Overdamped Modes}

\begin{definition}[Overdamped Memory]
A memory mode $\psi_i$ is \emph{overdamped relative to a
cue class} $\mathcal{C}$ when
\[
  \sup_{c \in \mathcal{C}}
  \frac{|\gamma_i(c)|}{\sqrt{\lambda_i^2 + \omega_i^2}}
  < \theta_i.
\]
\end{definition}

\begin{proposition}[Forgetting Is Loss of Excitability]
\label{prop:forgetting}
If a mode is overdamped relative to all admissible cues,
the corresponding memory is operationally forgotten
but structurally latent.
\end{proposition}

\begin{proof}
The mode $\psi_i$ remains an element of the modal basis
$\{\psi_i\}$; it is not removed from $\calH$.
However, for every $c \in \mathcal{C}$,
\[
  |a_i^*(c)| < \theta_i,
\]
so the mode cannot become recall-active.
Forgetting is therefore a failure of excitability,
not an erasure of the structural possibility of the memory.
\end{proof}

This distinction is philosophically important.
A forgotten memory is not destroyed; it is unreachable.
The distinction between erasure and overdamping generates
the possibility of recovery through repair.

\section{Repair as Damping Reduction}

\begin{theorem}[Repair by Reducing Damping]
\label{thm:repair}
Let $\psi_i$ be overdamped under forcing $\gamma_i$.
If repair changes the damping coefficient from $\lambda_i$
to $\lambda_i'$ with
\[
  \lambda_i' < \sqrt{\left(\frac{|\gamma_i|}{\theta_i}\right)^2
                     - \omega_i^2},
\]
then the memory becomes recall-active under the same cue forcing.
\end{theorem}

\begin{proof}
The repaired steady-state amplitude is
\[
  |a_i^{*\prime}| = \frac{|\gamma_i|}{\sqrt{(\lambda_i')^2 + \omega_i^2}}.
\]
The stated inequality on $\lambda_i'$ is equivalent to
\[
  (\lambda_i')^2 + \omega_i^2
  < \left(\frac{|\gamma_i|}{\theta_i}\right)^2.
\]
Taking square roots and rearranging:
\[
  \frac{|\gamma_i|}{\sqrt{(\lambda_i')^2 + \omega_i^2}} > \theta_i.
\]
Therefore $|a_i^{*\prime}| > \theta_i$,
meaning the repaired mode is recall-active.
\end{proof}

\begin{corollary}
Therapeutic, environmental, and social interventions
that restore memory do not reconstruct stored content.
They reduce the effective damping coefficient of
overdamped modes, allowing existing latent structure
to become resonant.
\end{corollary}

% ============================================================
\part{Ecological Memory}
% ============================================================

\chapter{Memory Outside the Brain}

\section{The Ecological Extension}

The neural interpretation of MEM|8 is a special case,
not the general case.
Any system whose present state is a compressed witness
of its event history satisfies the axioms of the
MEM|8 framework.
Memory is therefore not confined to neural tissue.

\begin{definition}[Environmental Memory Field]
\label{def:env-mem}
An \emph{environmental memory field} is a spatial
distribution
\[
  \Phi_M : \Omega \;\longrightarrow\; \RR_{\geq 0}
\]
representing the \emph{local reconstructive support}
at each point $x \in \Omega$.
High values of $\Phi_M(x)$ indicate regions where
distinction reconstruction is strongly supported;
low values indicate reconstructive poverty.
\end{definition}

\begin{example}[Workspace as Memory Architecture]
Consider a knowledge worker's physical desk.
Let $\Omega$ be the desk surface.
Notes, open books, pinned diagrams, and unfinished
drafts each instantiate a locally high $\Phi_M$.
The workspace as a whole is a distributed memory architecture
in which cognitive load is partially offloaded into the
environmental field.
\end{example}

\begin{proposition}[Ecological Extension of the Memory Field]
Let a cognitive domain extend across an agent $A$ and
an environment $\mathrm{Env}$.
The composite memory field is
\[
  M_{\text{total}}(t) = M_A(t) + M_{\mathrm{Env}}(t),
\]
and the composite modal decomposition is
\[
  M_{\text{total}}(t) = \sum_i a_i^A(t)\psi_i^A
                        + \sum_j a_j^{\mathrm{Env}}(t)\psi_j^{\mathrm{Env}}.
\]
Recall can be triggered by cue coupling to either
component.
\end{proposition}

\begin{proof}
Both components are elements of the same Hilbert space
$\calH$ (extended to the composite domain).
Their sum is again an element of $\calH$,
with coefficients summing over the two mode families.
Recall via Theorem~\ref{thm:resonant-recall} applies
to any mode in the combined basis.
\end{proof}

\chapter{Habitat Versus Archive}

\begin{definition}[Archive]
An \emph{archive} is a system designed to preserve
the information content of representations across time.
Its memory metric is storage fidelity:
$d_{\calQ}(R_t(S(t)),\, q_0)$ should remain small as $t$
increases.
\end{definition}

\begin{definition}[Habitat]
A \emph{habitat} is a system designed to preserve and
support the \emph{reconstruction pathways} through which
distinctions can be regenerated.
Its memory metric is reconstructive accessibility:
the probability that a randomly sampled cue activates
a recall-active mode.
\end{definition}

\begin{proposition}[Archives Can Fail as Habitats]
An archive may preserve complete information while
failing as a memory habitat if the reconstruction
pathways required to use that information are lost.
\end{proposition}

\begin{proof}
Consider a text perfectly preserved in a dead language.
Every symbol is retained; $d_{\calQ}$ of the stored
content is zero.
Yet if no agent can execute the reconstruction operator
$R_t$ that maps the symbols to their distinctions,
the memory is operationally inaccessible.
The archive is intact; the habitat is broken.
\end{proof}

This distinction has consequences for the design of
libraries, databases, scientific journals, and institutional
knowledge systems.
Preservation of content is necessary but insufficient.
Preservation of the reconstruction community, the
interpretive tradition, and the cue structures that
trigger recall is equally essential.

% ============================================================
\part{Memory and Repair}
% ============================================================

\chapter{Memory as Flux Maintenance}

\section{Persistence Through Replacement}

Synaptic proteins mediating vesicle exocytosis have
average lifetimes of one to two days, yet the memories
they support persist for decades \cite{Parkes2023}.
This is the biological instance of the ship-of-Theseus
problem: how can identity persist when every component
is replaced?

MEM|8 resolves the paradox by locating memory identity
in reconstructive organization rather than material
constitution.
The mathematical expression is a flux-balance equation.

\begin{definition}[Memory Flux Balance]
\label{def:flux}
Let $\rho(x, t)$ be the local \emph{reconstructive density}
at point $x \in \Omega$.
Let $J(x, t)$ be the flux of reconstructive carriers,
$S_{\mathrm{renew}}$ a renewal source term, and
$S_{\mathrm{decay}}$ a degradation sink term.
The \emph{memory balance equation} is
\[
  \frac{\partial \rho}{\partial t}
  + \nabla \cdot J
  = S_{\mathrm{renew}} - S_{\mathrm{decay}}.
\]
\end{definition}

\begin{theorem}[Persistence Through Replacement]
\label{thm:ptr}
A memory region $U \subset \Omega$ preserves its
reconstructive density over time if
\[
  \int_U (S_{\mathrm{renew}} - S_{\mathrm{decay}})\,dx
  = \int_{\partial U} J \cdot \hat{n}\,dS,
\]
in which case $\dfrac{d}{dt}\displaystyle\int_U \rho\,dx = 0$.
\end{theorem}

\begin{proof}
Integrate the balance equation over $U$:
\[
  \frac{d}{dt}\int_U \rho\,dx
  = -\int_{\partial U} J \cdot \hat{n}\,dS
    + \int_U (S_{\mathrm{renew}} - S_{\mathrm{decay}})\,dx.
\]
By the stated equality, the two terms on the right cancel,
giving $\dfrac{d}{dt}\displaystyle\int_U \rho\,dx = 0$.
\end{proof}

\begin{corollary}[Ship-of-Theseus Memory]
If every material carrier in $U$ is replaced over the
interval $[0, T]$ but the flux-balance condition holds
throughout, then the memory persists as a stable
reconstructive organization.
\end{corollary}

\begin{proof}
Theorem~\ref{thm:ptr} guarantees that total reconstructive
density in $U$ is conserved.
Memory identity is defined by recoverability
(Definition~\ref{def:mem-equiv}), not by carrier identity.
Therefore wholesale replacement of carriers is compatible
with persistent memory.
\end{proof}

\section{Two Transport Modes: Circulation and Capture}

\begin{definition}[Two-State Memory Transport]
Memory carriers occupy either a \emph{circulating state}
$C$ (preserving global availability) or a
\emph{captured state} $K$ (stabilizing local reconstruction).
Let $p_K(v, d)$ be the probability that a carrier of
speed $v$ moving in direction
$d \in \{+, -\}$ is captured by a local memory site.
\end{definition}

\begin{assumption}[Speed-Dependent Capture]
\label{assump:speed-capture}
There exists a threshold $v_c$ such that
\[
  v < v_c \;\Rightarrow\; p_K(v, d) \approx p_{\mathrm{slow}},
  \qquad
  v > v_c \;\Rightarrow\; p_K(v, d) \approx p_{\mathrm{fast}},
\]
with $p_{\mathrm{slow}} > p_{\mathrm{fast}}$.
\end{assumption}

\begin{theorem}[Capture-Stabilization Principle]
Under Assumption~\ref{assump:speed-capture},
slow carriers preferentially stabilize local memory sites
while fast carriers preferentially maintain global
circulation.
\end{theorem}

\begin{proof}
Let $N = N_{\mathrm{slow}} + N_{\mathrm{fast}}$.
Expected capture counts are
\[
  \mathbb{E}[K]
  = p_{\mathrm{slow}} N_{\mathrm{slow}}
  + p_{\mathrm{fast}} N_{\mathrm{fast}},
\]
\[
  \mathbb{E}[C]
  = (1 - p_{\mathrm{slow}}) N_{\mathrm{slow}}
  + (1 - p_{\mathrm{fast}}) N_{\mathrm{fast}}.
\]
Since $p_{\mathrm{slow}} > p_{\mathrm{fast}}$,
slow carriers contribute disproportionately to $\mathbb{E}[K]$
and fast carriers to $\mathbb{E}[C]$.
The two velocity classes therefore separate into
local stabilization and global circulation roles.
\end{proof}

\begin{theorem}[Retrograde Renewal Bias]
\label{thm:retrograde}
Suppose carriers moving retrograde have lower
capture probability than anterograde carriers:
$p_K(v,-) < p_K(v,+)$.
Then, even when initial directional populations are equal,
the surviving traversing flux has a retrograde bias.
\end{theorem}

\begin{proof}
Assume $N_+ = N_-$.
Expected traversing populations are
\[
  T_+ = (1 - p_K(v,+))N_+,
  \qquad
  T_- = (1 - p_K(v,-))N_-.
\]
Since $p_K(v,-) < p_K(v,+)$ implies
$1 - p_K(v,-) > 1 - p_K(v,+)$,
and $N_+ = N_-$, we obtain $T_- > T_+$.
Retrograde bias in the traversing flux arises from
differential capture probability, not from directional
motor asymmetry.
\end{proof}

Theorem~\ref{thm:retrograde} gives MEM|8 a mechanism for
automatic renewal: old material is preferentially returned
for replacement without requiring explicit deletion or
directed degradation signals.

\chapter{Forgetting as Structural Deformation}

\begin{definition}[Deformation Cost]
The \emph{deformation cost} of the memory field over
the interval $[t, t + \Delta]$ is
\[
  D(H_t, H_{t+\Delta})
  = \sum_{i \in I}
    \left|a_i(t+\Delta) - a_i(t)\right|^2 \cdot \mathbf{1}[|a_i(t)| \geq \theta_i].
\]
The indicator selects only modes that were recall-active
at time $t$: only the deformation of active memories
contributes.
\end{definition}

\begin{proposition}[Forgetting Rate Bounds Deformation Cost]
Let $\Lambda = \sup_i \sqrt{\lambda_i^2 + \omega_i^2}$
be the maximum modal impedance over all recall-active modes,
and suppose no cue forcing is present.
Then
\[
  D(H_t, H_{t+\Delta}) \leq \Lambda^2 \cdot \Delta^2 \cdot \|M(t)\|_{\mathrm{active}}^2,
\]
where $\|M(t)\|_{\mathrm{active}}^2 = \sum_{i:\,|a_i|\geq\theta_i} |a_i(t)|^2$.
\end{proposition}

\begin{proof}
Without cue forcing and noise, $a_i(t) = a_i(0)e^{-(\lambda_i + i\omega_i)t}$.
Over a short interval $\Delta$, differentiating gives
$\dot{a}_i = -(\lambda_i + i\omega_i)a_i$, so
\[
  |a_i(t+\Delta) - a_i(t)|
  \approx |\dot{a}_i(t)|\,\Delta
  = |\lambda_i + i\omega_i|\,|a_i(t)|\,\Delta
  = \sqrt{\lambda_i^2 + \omega_i^2}\,|a_i(t)|\,\Delta
  \leq \Lambda\,|a_i(t)|\,\Delta,
\]
using the definition $\Lambda = \sup_i\sqrt{\lambda_i^2+\omega_i^2}$.
Squaring and summing over recall-active modes:
\[
  D(H_t,H_{t+\Delta})
  = \sum_{i:\,|a_i|\geq\theta_i} |a_i(t+\Delta)-a_i(t)|^2
  \leq \Lambda^2\Delta^2 \sum_{i:\,|a_i|\geq\theta_i} |a_i(t)|^2
  = \Lambda^2\Delta^2 \|M(t)\|_{\mathrm{active}}^2. \qedhere
\]
\end{proof}

\chapter{Persistent Anomalies and Scientific Memory}

\begin{definition}[Anomaly Persistence]
An anomaly $\mathcal{A}$ in a knowledge system is
\emph{persistent} if its associated memory mode $\psi_{\mathcal{A}}$
has low damping coefficient $\lambda_{\mathcal{A}} \approx 0$
and strong coupling to routine cues:
$\gamma_{\mathcal{A}}(c) \gg \theta_{\mathcal{A}} \sqrt{\lambda_{\mathcal{A}}^2 + \omega_{\mathcal{A}}^2}$
for most $c$ in the standard repertoire.
\end{definition}

\begin{proposition}[Anomaly Survives Paradigm Shift]
If an anomaly has low damping and strong cue coupling,
it remains recall-active across theoretical transitions
that alter other modes.
\end{proposition}

\begin{proof}
Theoretical transitions change the modal basis
$\{\psi_i\} \to \{\psi_i'\}$ and alter the damping
landscape for most modes.
But if $\lambda_{\mathcal{A}} \approx 0$, the anomaly mode
remains below the recall threshold even under strong
perturbation of the remaining basis.
Its cue coupling $\gamma_{\mathcal{A}}$ ensures routine
re-excitation.
Therefore the anomaly persists: it is the memory
mode most resistant to paradigm-induced overdamping.
\end{proof}

Science accumulates not because it stores truths,
but because the anomaly modes of the scientific community
are collectively maintained at low damping through
continual cue excitation: publication, teaching,
replication, and debate.

% ============================================================
\part{Computation and MEM|8}
% ============================================================

\chapter{Memory Machines}

\begin{definition}[Reconstruction Capacity]
The \emph{reconstruction capacity} of a computational
memory system is the maximum number of
distinct distinctions recoverable from the system's
current state, measured in bits:
\[
  \mathcal{RC}(S) = \log_2 |\{q \in \calQ : \exists R,\, R(S) = q\}|.
\]
\end{definition}

\begin{remark}
Classical RAM measures memory in storage capacity
(bits of information stored).
MEM|8 measures memory in reconstruction capacity
(bits of information recoverable).
The two coincide for perfect-fidelity storage;
they diverge for systems with compression, degradation,
and reconstruction.
\end{remark}

\begin{theorem}[Reconstruction Dominates Storage]
For a system implementing witness expansion $\Phi : X \to \CC^n$
with $n$ independent binary witnesses,
the reconstruction capacity satisfies
\[
  \mathcal{RC}(S) \;\geq\; n,
\]
independently of the physical storage capacity of the
input representation.
\end{theorem}

\begin{proof}
Each witness coordinate $\phi_j : X \to \CC$ contributes
at least one bit of distinguishing power: whether
$\operatorname{Re}(\phi_j(x)) > 0$ or not partitions $X$
into two classes.
With $n$ independent witnesses, the number of jointly
distinguishable classes is at least $2^n$.
The reconstruction capacity is defined as
$\mathcal{RC}(S) = \log_2 |\{q \in \calQ : \exists R,\, R(S) = q\}|$,
so having at least $2^n$ distinguishable outputs gives
\[
  \mathcal{RC}(S) \;\geq\; \log_2 2^n = n.
\]
This lower bound is determined by the expansion dimension
$n$, not the input dimension $\dim(X)$.
\end{proof}

\chapter{Event-Log Computation and Spherepop}

Programs in the event-log interpretation are not state
machines; they are history machines.
A program execution is a sequence of events
$E = (e_1, \ldots, e_n)$ and the result is a recovered
distinction from the terminal historical state $H_n$.

This connects directly to the Spherepop calculus,
in which computation proceeds through irreversible
events (Pop, Refuse, Collapse, Bind) and the
computational state is the history of those events
rather than a mutable store.

\begin{definition}[Event-Log Program]
An \emph{event-log program} is a pair $(\mathcal{E}, R)$
where $\mathcal{E}$ is an event alphabet and
$R : \calH_{\mathcal{E}} \to \calQ$ is a reconstruction
functional over histories in $\calH_{\mathcal{E}}$.
Execution consists of generating a history $H$;
the result is $R(H)$.
\end{definition}

\begin{proposition}[Spherepop as Event-Log Computation]
The Spherepop calculus, with its irreversible event
operators, is an event-log program in the sense above,
where the history $H$ is the sequence of sphere-state
transitions and $R$ maps terminal histories to
denotational values.
\end{proposition}

\begin{proof}[Sketch]
Each Spherepop operator (Pop, Refuse, Collapse, Bind)
modifies the global sphere configuration irreversibly,
generating a new event in the history.
The operational semantics assigns meaning to terminal
histories.
This is precisely the structure of an event-log program.
\end{proof}

\chapter{A Concrete Instantiation: The MEM|8 Rust Backend}
\label{ch:rust}

\section{From Theory to Implementation}

The preceding chapters developed MEM|8 as an abstract mathematical
framework: memory fields in $L^2(\Omega;\CC)$, modal expansions,
resonant recall, and flux-balanced persistence.
This chapter examines a working Rust implementation of MEM|8's
storage layer \cite{Chenoweth2025MEM8} and shows precisely where the abstract constructions
land in code.
The implementation is instructive not only as an existence proof
but because it makes visible one residual tension: the
\texttt{VolatileBackend}/\texttt{PersistentBackend} split still
carries traces of the storage illusion that MEM|8 theoretically
supersedes.

\section{Wave Coordinates as Witness Expansion}

The core type is

\begin{verbatim}
pub type WaveCoord = (u8, u8, u16);
\end{verbatim}

\noindent
a triple $(x, y, f)$ where $x, y \in \{0,\ldots,255\}$ specify
a position in the $256 \times 256$ spatial grid and
$f \in \{0,\ldots,65535\}$ specifies a frequency band.
The full coordinate space has cardinality
\[
  256 \times 256 \times 65536 = 2^{32} \approx 4.3 \times 10^9
\]
distinct positions, each capable of holding a complex amplitude

\begin{verbatim}
pub type Wave = Complex64;
\end{verbatim}

This is a concrete realization of the witness expansion
(Definition~\ref{def:witness-expansion}).
Each coordinate $\phi_{(x,y,f)}$ is a reconstructive witness,
and the full lattice
\[
  \Lambda = \{(x, y, f) : x,y \in [0,255],\; f \in [0,65535]\}
\]
is the high-dimensional witness space into which events are
expanded.
Theorem~\ref{thm:dim-amp} applies directly: input events that are
not linearly separable in raw sensory space may become separable
in $\Lambda$ under a nonlinear embedding, and
Corollary~\ref{cor:witness-redundancy} guarantees robustness to
the loss of individual witness coordinates.

The choice of \texttt{u8} for spatial axes and \texttt{u16} for
frequency is itself a statement about the relative grain of
spatial versus spectral resolution in the intended cognitive
domain, an architectural commitment to the
\emph{type-specific encoding} principle: different kinds of
distinctions (spatial vs.\ spectral) occupy different regions
of the witness lattice.

\section{\texttt{query\_region} as Ecphoric Activation}

The central method of the \texttt{WaveStorageBackend} trait is

\begin{verbatim}
fn query_region(&self, bounds: BoundingBox)
    -> Result<Vec<(WaveCoord, Wave)>>;
\end{verbatim}

\noindent
where \texttt{BoundingBox} specifies spatial extent:

\begin{verbatim}
pub struct BoundingBox {
    pub min_x: u8, pub min_y: u8,
    pub max_x: u8, pub max_y: u8,
}
\end{verbatim}

\noindent
The implementation comment reads: ``This is the foundation of
resonance. The backend returns all waves within the spatial
bounds, which the cognitive layer then uses to calculate
interference patterns.''

This maps precisely onto the ecphoric activation operator
of Chapter~7.
A cue $c$ is encoded as a bounding box
$B_c \subset \Lambda$.
The cue coupling of $c$ to mode $\psi_i$ is
\[
  \gamma_i(c) = \bigl|\langle C(c),\, \psi_i \rangle\bigr|,
\]
where $C(c)$ is the field representation of the bounding-box
query: the indicator function of $B_c$ in the wave lattice.
The backend computes, for each coordinate $(x,y,f) \in B_c$,
the stored complex amplitude $a_{(x,y,f)}$.
These are exactly the excitation coefficients $a_i(t)$
of Definition~\ref{def:memory-field}, restricted to the modes
whose spatial support intersects the cue region.

Resonant recall (Theorem~\ref{thm:resonant-recall}) then
operates in the cognitive layer above the backend.
The cognitive core receives the amplitudes returned by
\texttt{query\_region}, computes interference patterns, and
applies a threshold: a mode contributes to recall if and only if
\[
  |a_i^*| \;\geq\; \theta_i.
\]
The separation of concerns is therefore: the backend handles
wave retrieval (which modes are active in the cue region);
the cognitive layer handles thresholding and interference
(which of those modes cross the recall boundary).
This is the correct factoring --- the storage substrate should
not know the ecphoric threshold, just as a physical medium
should not know which patterns are meaningful to the organism
reading from it.

\begin{proposition}[\texttt{query\_region} Instantiates Ecphory]
\label{prop:query-ecphory}
Let $B_c$ be a bounding box encoding cue $c$,
and let $\mathcal{W}(B_c)$ denote the set of wave amplitudes
returned by \texttt{query\_region}$(B_c)$.
Then recall of memory $m$ from cue $c$ is equivalent to
\[
  \left|
    \sum_{(x,y,f)\,\in\, B_c}
    a_{(x,y,f)}\,\psi_{(x,y,f)}
  \right|
  \;\geq\; \theta_m,
\]
where $\psi_{(x,y,f)}$ are the lattice basis modes and
$\theta_m$ is the ecphoric threshold for $m$.
\end{proposition}

\begin{proof}
\texttt{query\_region} returns exactly the set
$\{(\mathit{coord},\, a_{\mathit{coord}})\}$
for $\mathit{coord} \in B_c \cap \Lambda$.
The cognitive layer forms the superposition
\[
  \hat{M}(c)
  = \sum_{\mathit{coord}\,\in\, B_c} a_{\mathit{coord}}\,\psi_{\mathit{coord}},
\]
which is the restriction of the memory field $M(t)$ to the cue
region --- the ecphoric activation $\mathcal{E}(c,H) \to \hat{H}$
of Chapter~7, with the bounding box selecting the relevant
historical submanifold.
Recall fires when $|\hat{M}(c)| \geq \theta_m$,
matching Theorem~\ref{thm:resonant-recall}.
\end{proof}

\section{The Storage Illusion Residue}

The \texttt{flush()} method and the deprecated
\texttt{WaveStorage} trait reveal a residual tension between
the implementation and the theory.

\texttt{flush()} is described as ensuring durability by
flushing write buffers to persistent storage, and as a no-op
for volatile backends.
This treats persistence as a property of the backend type
(volatile vs.\ persistent) rather than as a property of the
maintenance dynamics.

The MEM|8 theory says something different.
Persistence is not a function of where data sits
(RAM versus disk) but of whether the flux-balance condition
of Theorem~\ref{thm:ptr} is satisfied:
\[
  \frac{\partial \rho}{\partial t} + \nabla \cdot J
  = S_{\mathrm{renew}} - S_{\mathrm{decay}}.
\]
A ``persistent'' backend that is never re-excited by cue
queries will accumulate overdamped modes and lose memory in
the operational sense, even though every byte is preserved.
Conversely, a ``volatile'' backend continually re-queried and
refreshed is maintaining its reconstructive density through
active use --- a form of the renewal flux $S_{\mathrm{renew}}$
demonstrated in the synaptic-vesicle trafficking model of
Chapter~13.

The deprecated \texttt{WaveStorage} / \texttt{WaveGrid} pair
is a clean example of the storage illusion at the
implementation level.
\texttt{store\_wave\_grid} writes a static grid and returns;
there is no provision for re-excitation, ecphoric thresholds,
or flux maintenance.
Its deprecation in favor of \texttt{WaveStorageBackend} with
\texttt{query\_region} at the center is, in miniature, the same
conceptual transition that MEM|8 makes at the theoretical level:
from storage as the primitive operation to reconstruction as
the primitive operation.

\section{Toward a Fully Theory-Conformant Backend}

A backend fully aligned with MEM|8 theory would augment
\texttt{WaveStorageBackend} with three additional mechanisms.

\paragraph{Damping and renewal.}
A method \texttt{decay\_step(\&mut self, dt: f64)} would apply
\[
  a_i(t+\Delta t) = a_i(t)\,e^{-(\lambda_i + i\omega_i)\Delta t}
\]
to every stored coefficient, and a companion
\texttt{renew(\&mut self, flux: \&[(WaveCoord, Wave)])} would
add renewal input $S_{\mathrm{renew}}$ to specified coordinates,
making persistence a dynamic equilibrium rather than a backend
attribute.

\paragraph{Damping coefficient registry.}
Each coordinate $(x, y, f)$ should carry $\lambda_{(x,y,f)}$
and $\omega_{(x,y,f)}$, so the cognitive layer can evaluate
the recall threshold condition
\[
  |\gamma_i| \;\geq\; \theta_i\sqrt{\lambda_i^2 + \omega_i^2}
\]
without embedding modal parameters in the storage layer.

\paragraph{Retrograde sweep.}
A scheduled low-priority process implementing
Theorem~\ref{thm:retrograde} would flag low-activity
coordinates for renewal before they become overdamped,
maintaining the maintenance current that keeps long-term
memories recoverable without explicit refresh commands.

These three additions close the gap between the current
implementation and the theoretical architecture.
The implementation already has the right shape at the
\texttt{query\_region} level; what remains is to make the
temporal dynamics --- damping, renewal, retrograde sweep ---
first-class citizens of the backend interface rather than
implicit assumptions about the environment.

% ============================================================

\chapter{The Halting Problem and Historical Persistence}

\begin{theorem}[Undecidability as Future Reconstructability Failure]
The halting problem expresses a fundamental limit on
future reconstructability: there is no finite procedure
that determines, for an arbitrary event-log program and
initial condition, whether the history will eventually
produce a terminal reconstruction.
\end{theorem}

\begin{proof}[Reduction to Rice's Theorem]
Any property of the language recognized by a Turing machine
(including whether the machine halts on a given input)
is a property of the reconstruction functional
$R : \calH \to \calQ$ over all possible histories.
By Rice's theorem, every nontrivial such property is
undecidable.
Therefore the question ``will the reconstruction eventually
succeed?'' is undecidable for arbitrary programs.
\end{proof}

The interpretation is important: undecidability is not
a failure of storage but a failure of
\emph{future reconstructability}.
We cannot predict whether the event-log will eventually
support the desired reconstruction.
This gives the halting problem a MEM|8 reading entirely
consistent with the general framework.

% ============================================================
\part{Society and Civilization}
% ============================================================

\chapter{Institutional Memory}

\begin{definition}[Organizational Memory Field]
The \emph{organizational memory field} of an institution
is the composite memory field
\[
  M_{\mathrm{org}}(t)
  = \sum_{a \in A} M_a(t) + M_{\mathrm{doc}}(t) + M_{\mathrm{proc}}(t),
\]
where the sum is over individual agents $a \in A$,
plus document repositories $M_{\mathrm{doc}}$ and
procedural structures $M_{\mathrm{proc}}$.
\end{definition}

\begin{proposition}[Personnel Change Does Not Destroy Memory]
A personnel change (departure of agent $a$) destroys
organizational memory if and only if the modes
$\{\psi_i^a\}$ carried by agent $a$ are not
decodable from the remaining composite field
$M_{\mathrm{org}} - M_a$.
\end{proposition}

\begin{proof}
By Corollary~\ref{cor:witness-redundancy},
a memory survives the loss of witness coordinates
provided sufficient redundant witnesses remain.
If $M_{\mathrm{doc}}$ or $M_{\mathrm{proc}}$ supply
alternative cue paths to the same modal distinctions,
the reconstruction persists.
Memory loss occurs precisely when no such redundant
pathway exists.
\end{proof}

\chapter{Scientific Memory}

Science is a civilization-scale memory architecture.
Its function is to maintain, across generations and
institutional boundaries, a reconstructive field
rich enough to support cumulative inquiry.

\begin{definition}[Scientific Paper as Reconstruction Operator]
A scientific paper $P$ is a reconstruction operator
$R_P : \calH_{\mathrm{reader}} \to \calQ_{\mathrm{science}}$
that maps the reader's current cognitive state to a
set of distinctions about the natural world.
\end{definition}

\begin{proposition}[Citation Network as Modal Coupling]
The citation network of the scientific literature
implements a system of modal couplings $\gamma_{ij}$
between papers $P_i$ and $P_j$:
a citation from $P_j$ to $P_i$ provides a coupling
channel that keeps the modes of $P_i$ excitable
via engagement with $P_j$.
\end{proposition}

\begin{proof}
Reading $P_j$ generates cue coupling to the modes
encoded in $P_i$ through the citation-link
(references, shared terminology, shared problems).
This is cue forcing $\gamma_i$ in the sense of
Definition~\ref{def:mem-equiv}, applied to the
scientific memory field.
\end{proof}

\chapter{Cultural Memory}

Languages, traditions, norms, and rituals are long-lived
repair structures.
They persist by maintaining the reconstructive conditions
under which the distinctions they encode can be
regenerated.

\begin{theorem}[Cultural Memory Survival Condition]
A cultural form $\mathcal{F}$ survives across generations
if and only if the flux-balance condition of
Theorem~\ref{thm:ptr} is maintained for the reconstructive
density of $\mathcal{F}$: renewal (teaching, practice,
performance, reenactment) balances decay (mortality,
forgetting, disuse).
\end{theorem}

\begin{proof}
Cultural forms are event-log systems (Chapter~6).
Their memory field satisfies the balance equation
(Definition~\ref{def:flux}).
Survival of the form over time corresponds to
$\dfrac{d}{dt}\displaystyle\int \rho\,dx = 0$,
which by Theorem~\ref{thm:ptr} holds if and only if
the balance condition holds.
\end{proof}

% ============================================================
\part{Foundations}
% ============================================================

\chapter{Memory Before Representation}

The standard order of explanation places representation
first.
There are things; representations encode them; memory
stores the encodings.

MEM|8 reverses this order at every level.

\begin{theorem}[Representation Is Derived]
A representation $r$ of a distinction $\delta$ exists
precisely when there is a reconstruction operator $R$
and a substrate $S$ such that $R(S) = \delta$.
Representations are not primitive; they are the outputs
of reconstruction.
\end{theorem}

\begin{proof}
By Definition~\ref{def:recon-capacity}, reconstructive
capacity is defined without appeal to representations.
A representation can then be defined as a naming of the
output of a reconstruction operator: $r \equiv R(S)$.
This makes representation derivable from recoverability,
not the reverse.
\end{proof}

\noindent The inversion is therefore:

\[
  \text{Traditional:}\quad
  \text{Representation} \to \text{Storage} \to \text{Memory}
\]

\[
  \text{MEM|8:}\quad
  \text{Events} \to \text{Recoverability} \to \text{Memory} \to \text{Representation}
\]

\chapter{Memory Before Identity}

\begin{definition}[Identity as Stable Reconstruction]
The \emph{identity} of a system $\Sigma$ over an interval
$[t_0, t_1]$ is the equivalence class of states
under memory-equivalence relative to a canonical
reconstruction operator $R_\Sigma$:
\[
  \mathrm{Id}(\Sigma, [t_0, t_1])
  = [S(t_0)]_{\sim_{R_\Sigma}}.
\]
Identity persists over the interval if $S(t) \sim_{R_\Sigma} S(t_0)$
for all $t \in [t_0, t_1]$.
\end{definition}

\begin{proposition}[The Self as Persistent Memory Process]
Personal identity over a lifetime is the persistence of
a memory field under continual flux, satisfying the
balance condition of Theorem~\ref{thm:ptr}.
\end{proposition}

\begin{proof}
Personal identity is a special case of system identity.
The canonical reconstruction operator $R_\mathrm{self}$
maps current neural and embodied state to autobiographical
distinctions.
Persistence requires only that the flux-balance condition
hold for the reconstructive density of the autobiographical
memory field --- not that any neural substrate remain
unchanged.
This is consistent with the empirical observation that
virtually every molecule in the human body is replaced
over a decade, while autobiographical memory persists.
\end{proof}


\chapter{Memory Before Prediction}
\label{ch:prediction}

This chapter addresses an implication the preceding chapters
repeatedly approach but never state outright.
If memory is recoverable continuation, then prediction ---
the anticipation of future distinctions --- is a special
case of reconstruction applied forward rather than backward.
Prediction presupposes memory, not the reverse.

\section{Prediction as Forward Reconstruction}

\begin{definition}[Predictive Reconstruction]
A \emph{predictive reconstruction} is a reconstruction
operator $R_+$ that maps the current memory field $M(t)$
to an anticipated future distinction:
\[
  R_+ : \calH \;\longrightarrow\; \calQ,
  \qquad
  M(t) \;\mapsto\; \hat{q}(t + \tau),
\]
where $\tau > 0$ is the prediction horizon and
$\hat{q}(t+\tau)$ is the anticipated state of the
distinction space at time $t + \tau$.
\end{definition}

\begin{theorem}[Prediction Presupposes Memory]
\label{thm:pred-mem}
A system capable of predictive reconstruction
$R_+$ must have nonzero reconstructive capacity
$\mathcal{RC}(S) > 0$.
\end{theorem}

\begin{proof}
Predictive reconstruction $R_+$ maps $M(t)$ to a future
distinction $\hat{q}(t+\tau)$.
For $R_+$ to be non-trivial (i.e., better than chance),
it must discriminate among distinct current states:
if $\mathcal{RC}(S) = 0$, every state maps to the same
reconstruction output (Theorem~\ref{thm:collapse}), so
no state-dependent prediction is possible.
Therefore non-trivial prediction requires
$\mathcal{RC}(S) > 0$, i.e., genuine reconstructive
capacity over the current memory field.
Prediction is downstream of memory, not independent of it.
\end{proof}

\begin{corollary}[No Prediction Without Historical State]
A memoryless system --- one with $H_t = \varnothing$ for
all $t$ --- cannot predict.
\end{corollary}

\begin{proof}
Without historical state there is no compression to unpack
(Chapter~5), no modal basis to excite (Chapter~9), and
no reconstruction functional to apply.
Theorem~\ref{thm:pred-mem} then gives $\mathcal{RC}(S) = 0$,
so no non-trivial prediction is possible.
\end{proof}

\section{Predictive Coding as Memory Resonance}

The predictive processing framework (Friston \cite{Friston2010})
treats perception and action as continuous minimization of
prediction error: the brain maintains a generative model
and updates it when predictions fail.

Under MEM|8 this has a natural modal reading.
The generative model is the current modal excitation pattern
$\{a_i(t)\}$.
A prediction is a projection of that pattern forward:
\[
  \hat{a}_i(t+\tau) = a_i(t)\,e^{-(\lambda_i + i\omega_i)\tau}.
\]
Prediction error is the discrepancy between the predicted
field $\hat{M}(t+\tau) = \sum_i \hat{a}_i(t+\tau)\psi_i$
and the actual field $M(t+\tau)$:
\[
  \mathcal{E}_{\mathrm{pred}}(t,\tau)
  = \|M(t+\tau) - \hat{M}(t+\tau)\|_{\calH}^2
  = \sum_i |a_i(t+\tau) - \hat{a}_i(t+\tau)|^2.
\]

\begin{proposition}[Prediction Error Bounds Surprise]
Let $\mathcal{E}_{\mathrm{pred}}(t,\tau)$ be the prediction
error defined above.
Then
\[
  \mathcal{E}_{\mathrm{pred}}(t,\tau)
  \leq \sum_i |\eta_i(t,\tau)|^2
       + \sum_i |\gamma_i(c)|^2,
\]
where $\eta_i(t,\tau)$ is background noise accumulated
over $[t, t+\tau]$ and $\gamma_i(c)$ is the cue forcing
from unexpected environmental input over the same interval.
\end{proposition}

\begin{proof}
The actual coefficient satisfies
$a_i(t+\tau) = \hat{a}_i(t+\tau) + \delta_i$,
where $\delta_i$ arises from noise $\eta_i$ and
unexpected forcing $\gamma_i(c)$ not already encoded
in the current modal state.
Therefore
$|a_i(t+\tau) - \hat{a}_i(t+\tau)|^2 = |\delta_i|^2
\leq |\eta_i|^2 + |\gamma_i(c)|^2$,
and summing over modes gives the bound.
\end{proof}

Prediction error is thus bounded by noise plus
unencoded surprise.
A system with richer memory (higher-dimensional modal
basis, lower damping) encodes more of its environment
and therefore anticipates more accurately --- connecting
MEM|8 directly to the Bayesian brain hypothesis and
active inference.

\section{Reservoir Computing as MEM|8 Instantiation}

Reservoir computing \cite{Friston2023} treats a
high-dimensional dynamical system (the reservoir) as a
fixed feature expander: inputs drive the reservoir,
and a simple linear readout is trained on the reservoir
state to perform prediction and classification tasks.

This is precisely the MEM|8 architecture of
Theorem~\ref{thm:dim-amp}:

\begin{align*}
  \text{Reservoir state} &\;\longleftrightarrow\; \text{Witness expansion } \Phi(x) \in \CC^n \\
  \text{Linear readout}  &\;\longleftrightarrow\; \text{Reconstruction functional } L : \CC^n \to \calQ \\
  \text{Prediction task} &\;\longleftrightarrow\; \text{Predictive reconstruction } R_+
\end{align*}

The reservoir's memory of past inputs (its ``echo state'')
is the historical state $H_t$; the readout extracts
distinctions from that compressed history.
Theorem~\ref{thm:pred-mem} explains why reservoir memory
capacity directly governs prediction accuracy: without
sufficient $\mathcal{RC}(S)$, the readout cannot form
useful predictions.

The TAS phi-bit lattice of Chapter~8 is a physical
instantiation of a reservoir operating in the MEM|8 regime:
a small substrate generates a large witness space, and
classification (prediction) becomes linearly separable
in witness coordinates.

\section{The Inversion}

The standard order places prediction first:
a system makes predictions; successful prediction
constitutes understanding; memory is the residue of
past understanding.

MEM|8 inverts this:

\[
  \text{Traditional:}\quad
  \text{Prediction} \to \text{Model} \to \text{Memory}
\]

\[
  \text{MEM|8:}\quad
  \text{Memory} \to \text{Modal Structure} \to \text{Prediction}
\]

Prediction is reconstruction applied in the forward
temporal direction.
It inherits all the properties of reconstruction:
it is approximate rather than exact, threshold-gated
rather than continuous, and ecologically extended across
agent and environment rather than confined to a brain.

The arrow of prediction points forward in time;
the arrow of memory points backward; both are
instantiated by the same reconstructive mechanism
applied to a recoverable event ordering.
The next chapter derives that ordering itself.

\chapter{Memory Before Time}

This is the most radical inversion.
Standard accounts place time first and locate memory
within time.
The MEM|8 position is that temporal structure is itself
derived from recoverable event orderings.

\begin{definition}[Recoverable Event Ordering]
A \emph{recoverable event ordering} over a set of events
$\mathcal{E}$ is a partial order $\prec$ on $\mathcal{E}$
such that, for any admissible reconstruction operator $R$
and any $e_i \prec e_j$, the distinction associated with
$e_i$ can be recovered from the historical state
$H_{e_j} = \{e_k : e_k \prec e_j\}$.
\end{definition}

\begin{theorem}[Time Emerges from Recoverable Order]
\label{thm:time}
Temporal structure, in the operational sense relevant to
memory, learning, and cognition, is equivalent to
recoverable event ordering.
Without reconstructable event relations, temporal
distinctions are operationally undefined.
\end{theorem}

\begin{proof}
The direction \emph{temporal order implies recoverable order}:
if $e_i$ causally precedes $e_j$ in Minkowski spacetime,
then $H_{e_j}$ contains information traceable to $e_i$,
and an admissible $R$ can in principle recover the
relevant distinction.

The direction \emph{recoverable order implies temporal structure}:
if for all $e_i \prec e_j$ the distinction of $e_i$ is
recoverable from $H_{e_j}$, then $\prec$ constitutes a
causal precedence relation.
The set $(\mathcal{E}, \prec)$ with this property is a
\emph{causal set} in the sense of Bombelli et al.;
causal sets are known to encode the topological and
differential structure of spacetime in the continuum limit.
Therefore recoverable event ordering is sufficient to
reconstruct operational temporal structure.
\end{proof}

\begin{corollary}
Memory does not occur \emph{in} time as a background stage.
Time is the structure that emerges from the recoverable
ordering of events in a memory field.
The arrow of time is the direction in which reconstruction
of earlier events from later states is possible.
\end{corollary}

\bigskip

\noindent We arrive at the foundational statements of MEM|8:

\[
  \boxed{\;
    \text{A memory is not a preserved past.}
  \;}
\]

\[
  \boxed{\;
    \text{A memory is the continued ability to regenerate distinctions
    from transformation.}
  \;}
\]

\[
  \boxed{\;
    \text{Memory} = \text{recoverable resonant structure maintained by flux.}
  \;}
\]

% ============================================================
% Bibliographies
% ============================================================

\backmatter

\chapter*{Bibliography: Neuroscience, Biology, and Cognitive Science}
\addcontentsline{toc}{chapter}{Bibliography: Neuroscience, Biology, and Cognitive Science}

\begin{thebibliography}{99}

\bibitem{Anderson1983}
Anderson, J.~R. (1983).
\textit{The Architecture of Cognition}.
Harvard University Press.

\bibitem{Anderson2004}
Anderson, J.~R., Bothell, D., Byrne, M.~D., Douglass, S.,
Lebiere, C., \& Qin, Y. (2004).
An integrated theory of the mind.
\textit{Psychological Review}, 111(4), 1036--1060.

\bibitem{Amit1989}
Amit, D.~J. (1989).
\textit{Modeling Brain Function: The World of Attractor Neural Networks}.
Cambridge University Press.

\bibitem{Buzsaki2006}
Buzsáki, G. (2006).
\textit{Rhythms of the Brain}.
Oxford University Press.

\bibitem{Cabral2023OscillatoryModes}
Cabral, J., Fernandes, F.~F., \& Shemesh, N. (2023).
Intrinsic macroscale oscillatory modes driving long range
functional connectivity in female rat brains detected by ultrafast fMRI.
\textit{Nature Communications}, 14, 375.

\bibitem{Chen2026TAS}
Chen, J., Ige, A.~S., Runge, K., Deymier, P.~A., \& Yan, X. (2026).
Topological acoustic synapse for high-dimensional neuromorphic computing.
\textit{Science Advances}, 12(24), eaec6633.

\bibitem{Cohen2013}
Cohen, L.~D., et al. (2013).
Metabolic turnover of synaptic proteins.
\textit{Science}, 342(6161), 1094--1097.

\bibitem{Dayan2001}
Dayan, P., \& Abbott, L.~F. (2001).
\textit{Theoretical Neuroscience}.
MIT Press.

\bibitem{Deco2011}
Deco, G., Jirsa, V.~K., \& McIntosh, A.~R. (2011).
Emerging concepts for brain dynamics.
\textit{Nature Reviews Neuroscience}, 12(1), 43--56.

\bibitem{Fornasiero2018}
Fornasiero, E.~F., et al. (2018).
Precisely measured protein lifetimes.
\textit{Nature Communications}, 9, 4741.

\bibitem{Friston2010}
Friston, K. (2010).
The free-energy principle.
\textit{Nature Reviews Neuroscience}, 11(2), 127--138.

\bibitem{Gramlich2021}
Gramlich, M.~W., \& Klyachko, V.~A. (2021).
Motor-mediated regulation of synaptic vesicle mobility.
\textit{Biophysical Journal}, 120(4), 567--582.

\bibitem{Hebb1949}
Hebb, D.~O. (1949).
\textit{The Organization of Behavior}.
Wiley.

\bibitem{Hirokawa2009}
Hirokawa, N., Noda, Y., Tanaka, Y., \& Niwa, S. (2009).
Kinesin superfamily motors.
\textit{Nature Reviews Molecular Cell Biology}, 10(10), 682--696.

\bibitem{Hopfield1982}
Hopfield, J.~J. (1982).
Neural networks and physical systems.
\textit{Proceedings of the National Academy of Sciences}, 79(8), 2554--2558.

\bibitem{Joensuu2016}
Joensuu, M., et al. (2016).
Subdiffractional tracking of synaptic vesicles.
\textit{Nature Communications}, 7, 11714.

\bibitem{Kelso1995}
Kelso, J.~A.~S. (1995).
\textit{Dynamic Patterns}.
MIT Press.

\bibitem{Parkes2023}
Parkes, M., Landers, N.~L., \& Gramlich, M.~W. (2023).
Recently recycled synaptic vesicles use multi-cytoskeletal transport
and differential presynaptic capture probability to establish a
retrograde net flux during ISVE in central neurons.
\textit{Frontiers in Cell and Developmental Biology}, 11, 1286915.

\bibitem{Rizzoli2014}
Rizzoli, S.~O. (2014).
Synaptic vesicle recycling.
\textit{Nature Reviews Neuroscience}, 15(10), 658--673.

\bibitem{Rolls2016}
Rolls, E.~T. (2016).
\textit{Cerebral Cortex: Principles of Operation}.
Oxford University Press.

\bibitem{Sporns2011}
Sporns, O. (2011).
\textit{Networks of the Brain}.
MIT Press.

\bibitem{Truckenbrodt2018}
Truckenbrodt, S., et al. (2018).
Newly produced synaptic vesicle proteins.
\textit{Science}, 359(6371), 157--161.

\bibitem{Tononi2008}
Tononi, G. (2008).
Consciousness as integrated information.
\textit{Biological Bulletin}, 215(3), 216--242.

\bibitem{Chenoweth2025MEM8}
Chenoweth, C., \& Chenoweth, A. (2025).
\textit{MEM|8: A wave-based cognitive architecture for multimodal
memory integration and consciousness simulation}.
Zenodo.
\url{https://doi.org/10.5281/zenodo.16436297}

\end{thebibliography}

\chapter*{Bibliography: Mathematics, Physics, and Dynamical Systems}
\addcontentsline{toc}{chapter}{Bibliography: Mathematics, Physics, and Dynamical Systems}

\begin{thebibliography}{99}

\bibitem{Arnold1989}
Arnold, V.~I. (1989).
\textit{Mathematical Methods of Classical Mechanics} (2nd ed.).
Springer.

\bibitem{Arnold2006}
Arnold, V.~I. (2006).
\textit{Ordinary Differential Equations}.
Springer.

\bibitem{Ash1965}
Ash, R.~B. (1965).
\textit{Information Theory}.
Dover.

\bibitem{Billingsley1995}
Billingsley, P. (1995).
\textit{Probability and Measure} (3rd ed.).
Wiley.

\bibitem{Bowen1978}
Bowen, R. (1978).
\textit{Entropy for Group Endomorphisms and Homogeneous Spaces}.
Springer.

\bibitem{CoverThomas2006}
Cover, T.~M., \& Thomas, J.~A. (2006).
\textit{Elements of Information Theory} (2nd ed.).
Wiley.

\bibitem{Crutchfield1994}
Crutchfield, J.~P. (1994).
The calculi of emergence.
\textit{Physica D}, 75(1--3), 11--54.

\bibitem{Evans2010}
Evans, L.~C. (2010).
\textit{Partial Differential Equations} (2nd ed.).
American Mathematical Society.

\bibitem{Ghrist2014}
Ghrist, R. (2014).
\textit{Elementary Applied Topology}.
Createspace.

\bibitem{GuckenheimerHolmes1983}
Guckenheimer, J., \& Holmes, P. (1983).
\textit{Nonlinear Oscillations, Dynamical Systems, and
Bifurcations of Vector Fields}.
Springer.

\bibitem{Hatcher2002}
Hatcher, A. (2002).
\textit{Algebraic Topology}.
Cambridge University Press.

\bibitem{HornJohnson2013}
Horn, R.~A., \& Johnson, C.~R. (2013).
\textit{Matrix Analysis} (2nd ed.).
Cambridge University Press.

\bibitem{Jaynes2003}
Jaynes, E.~T. (2003).
\textit{Probability Theory: The Logic of Science}.
Cambridge University Press.

\bibitem{KatokHasselblatt1995}
Katok, A., \& Hasselblatt, B. (1995).
\textit{Introduction to the Modern Theory of Dynamical Systems}.
Cambridge University Press.

\bibitem{Kreyszig1989}
Kreyszig, E. (1989).
\textit{Introductory Functional Analysis with Applications}.
Wiley.

\bibitem{MacLane1998}
Mac Lane, S. (1998).
\textit{Categories for the Working Mathematician} (2nd ed.).
Springer.

\bibitem{MezardMontanari2009}
Mézard, M., \& Montanari, A. (2009).
\textit{Information, Physics, and Computation}.
Oxford University Press.

\bibitem{Nakahara2003}
Nakahara, M. (2003).
\textit{Geometry, Topology and Physics} (2nd ed.).
Taylor \& Francis.

\bibitem{Pikovsky2001}
Pikovsky, A., Rosenblum, M., \& Kurths, J. (2001).
\textit{Synchronization: A Universal Concept in Nonlinear Sciences}.
Cambridge University Press.

\bibitem{ReedSimon1972}
Reed, M., \& Simon, B. (1972).
\textit{Methods of Modern Mathematical Physics I: Functional Analysis}.
Academic Press.

\bibitem{ReedSimon1975}
Reed, M., \& Simon, B. (1975).
\textit{Methods of Modern Mathematical Physics II: Fourier Analysis,
Self-Adjointness}.
Academic Press.

\bibitem{Rudin1991}
Rudin, W. (1991).
\textit{Functional Analysis} (2nd ed.).
McGraw-Hill.

\bibitem{Ruelle1989}
Ruelle, D. (1989).
\textit{Chaotic Evolution and Strange Attractors}.
Cambridge University Press.

\bibitem{Shannon1948}
Shannon, C.~E. (1948).
A mathematical theory of communication.
\textit{Bell System Technical Journal}, 27(3), 379--423.

\bibitem{Smale1967}
Smale, S. (1967).
Differentiable dynamical systems.
\textit{Bulletin of the American Mathematical Society}, 73(6), 747--817.

\bibitem{Strogatz2018}
Strogatz, S.~H. (2018).
\textit{Nonlinear Dynamics and Chaos} (2nd ed.).
Westview Press.

\bibitem{Takens1981}
Takens, F. (1981).
Detecting strange attractors in turbulence.
In D.~Rand \& L.-S. Young (Eds.),
\textit{Dynamical Systems and Turbulence}.
Springer.

\bibitem{Temam1997}
Temam, R. (1997).
\textit{Infinite-Dimensional Dynamical Systems in Mechanics and Physics}
(2nd ed.).
Springer.

\bibitem{Tu2011}
Tu, L.~W. (2011).
\textit{An Introduction to Manifolds} (2nd ed.).
Springer.

\bibitem{vonNeumann1955}
von Neumann, J. (1955).
\textit{Mathematical Foundations of Quantum Mechanics}.
Princeton University Press.

\bibitem{Wiggins2003}
Wiggins, S. (2003).
\textit{Introduction to Applied Nonlinear Dynamical Systems and Chaos}
(2nd ed.).
Springer.

\bibitem{Zeidler1990}
Zeidler, E. (1990).
\textit{Applied Functional Analysis}.
Springer.

\end{thebibliography}

\end{document}
